%---------------------------------------------------------------------------%
%->> Backmatter
%---------------------------------------------------------------------------%
\chapter[致谢]{致\quad 谢}\chaptermark{致\quad 谢}% syntax: \chapter[目录]{标题}\chaptermark{页眉}
%\thispagestyle{noheaderstyle}% 如果需要移除当前页的页眉
%\pagestyle{noheaderstyle}% 如果需要移除整章的页眉

此处填写致谢。


\rightline{2023年6月}
\chapter{作者简历及攻读学位期间发表的学术论文与其他相关学术成果}

\section*{作者简历：}
2019年09月——2023年06月，在宁波大学信息科学与工程学院（系）获得工学学士学位。

% ××××年××月——××××年××月，在××大学××院（系）获得硕士学位。

2023年09月——2026年06月，在中国科学院信息工程研究所硕士学位。

% 工作经历：

\section*{已发表（或正式接受）的学术论文：}

{
\setlist[enumerate]{}% restore default behavior
\begin{enumerate}[nosep]
    \item Y. Lu et al., "A Fine-Grained Access Control for User Carbon Emission Data Based on Spatio-Temporal Attribute-Based Encryption," 2024 4th International Conference on Electronic Information Engineering and Computer Communication (EIECC), Wuhan, China, 2024, pp. 797-801, doi: 10.1109/EIECC64539.2024.10929490.

    \item Yansen Xin, Rui Zhang, Zhenglin Fan and Ze Jia. Dobby: A Private Time Series Data Analytics System with Enforcement of Flexible Policies
\end{enumerate}
}

% \section*{申请或已获得的专利：}

% （无专利时此项不必列出）

\section*{参加的研究项目及获奖情况：}
% \setlist[enumerate]
\begin{enumerate}[nosep]
    \item 南方电网区块链平台建设与维护
    \item 基于区块链的用电碳排放数据存储及系统构建关键技术研究
    \item 基于数据水印的立体化数据泄漏预防与溯源技术研究
    \item 数据强制性受控共享系统研发
\end{enumerate}
% restore default behavior

\cleardoublepage[plain]% 让文档总是结束于偶数页，可根据需要设定页眉页脚样式，如 [noheaderstyle]
%---------------------------------------------------------------------------%
